\section{Introduction}
\label{secIntroductionArchiReseau}

Ce document décrit l'architecture réseau \textbf{effectivement retenue} sur
l'installation réelle de la cellule robotique de la salle C180 : organes
physiques concernés, séparation des réseaux Ethernet/IP et adressage IP de
chaque équipement.

Pour la démarche de conception ayant conduit à ce choix (contraintes prises
en compte, séparation des réseaux, plan d'adressage), voir le document
\texttt{01\_conception}, en particulier la contrainte de séparation réseau
qui y est posée (\ref{secContraintes}).

Ce document renvoie par ailleurs vers d'autres documents de référence pour
les points qui y sont détaillés, afin de ne pas les redécrire :
\begin{itemize}
    \item le détail des \textbf{données échangées} avec chaque équipement
    (poste guichet, robot, IHM/pupitre) : voir le document
    \texttt{02b\_donnees\_echangees}, sections \ref{secDonneesGuichet},
    \ref{secDonneesRobot} et \ref{secDonneesIhm} ;
    \item la \textbf{procédure de configuration} des coupleurs réseau de
    l'automate M340 : voir le document \texttt{02c\_configuration\_noc},
    section \ref{secEtapeConfigReseau} ;
    \item les \textbf{préfixes de nommage} des variables associées au
    pupitre opérateur : voir le document \texttt{00\_conventions\_nommage},
    section \ref{secVariablesBooleennes}.
\end{itemize}
